\documentclass[12pt]{article}
\usepackage[shortlabels]{enumitem}
\usepackage{float}
\usepackage{xcolor}
\usepackage[a4paper, top=1.5cm, bottom=1.5cm, left=2cm, right=2cm]{geometry}
\usepackage{parskip}
\usepackage{setspace}
\usepackage{lmodern}
\usepackage[T1]{fontenc}
\usepackage[brazil]{babel}
\usepackage{hyperref}
\usepackage{graphicx}
\usepackage{amsmath}
\usepackage{amssymb}
\usepackage{amsfonts}
\usepackage{dsfont}
\usepackage{float}
\title{Lista 1 - Estatística para Administração - 2025.2}
\begin{document}
\date{}
\maketitle


1) Evidenciar o fato de que, em muitos casos, recursos como dinheiro e tempo são limitados. Outro fator é  a impossibilidade físicas,
 como é o caso do exame de sangue, em que a extração de todo o sangue implicaria na morte do paciente. 

2) Amostragem probabilística envolve o uso de métodos aleatórios para a seleção dos elementos da amostra, não havendo nenhuma interferência por parte do condutor da pesquisa na seleção da amostra. 
Já a amostragem não probabilística, não se vale da aleatoriedade para seleção dos elementos da amostra, possuindo interferência direta do condutor da pesquisa na seleção dos elementos que irão compor a amostra. 
Não é possível dizer que a amostragem probabilística é superior à não probabilística, entretanto, podemos dizer que a amostragem probabilística oferece grandes vantagens, como, por exemplo,
a possibilidade de realizar o processo de inferência estatística. 

3) Existem várias vantagens e desvantagens que discutimos em sala. Algumas vantagens são: método simples e de fácil execução, garante que cada indivíduo da população tenha a mesma chance de ser selecionado. 
Algumas desvantagen são: Exige conhecimento prévio acerca dos elementos que contsituem a população, além disso, pode incorrer na seleção de uma amostra concentrada em apenas uma parcela da população. 

4) Considere a seguinte situação:

\begin{enumerate}[a)]
    \item Alunos da escola na qual Leticia estudou.
    \item O fator principal foi o tempo, uma vez que Leticia não teria tempo hábil para entrevistar todos os alunos. 
    A escolha de trabalhar com a amostra é razoável, haja vista a impossibilidade de  entrevistar todos os alunos. 
    \item O uso da Amostragem Aleatória Estratificada garante que a representatividade de cada estrato na população seja refletida na amostra, evitando que algum grupo
    seja subamostrado ou ignorado. Por outro lado, exige conhecimento prévio acerca da representatividade de cada estrato na população, o que, no caso de Leticia, não é um problema, já que a diretora forneceu essa informação. 
    \item Estratos são grupos homogêneos em seu interior, mas heterogêneos entre si. Os Estratos são: Alunos do sexo masculino do Ensino Fundamental,
     Alunos do sexo feminino do Ensino Fundamental, Alunos do sexo masculino do Ensino Médio e Alunos do sexo feminino do Ensino Médio.
    \item  \begin{itemize}
        \item Alunos do sexo masculino do Ensino Fundamental: 19.
        \item Alunos do sexo feminino do Ensino Fundamental: 30.
        \item  Alunos do sexo masculino do Ensino Médio: 48.
        \item Alunos do sexo feminino do Ensino Médio: 43.
    \end{itemize}
    \item 35.
    \item \begin{itemize}
        \item Alunos do sexo feminino: 73
        \item Alunos do sexo masculino: 67
    \end{itemize} 
    \item A abordagem proporcional exige conhecimento prévio acerca da proporção de indivíduos em cada estrato, no entanto, garante que a proporção
     do estrato na população seja mantida na amostra, gerando uma amostra mais parecida com a população quanto à característica analisada. 
    A abordagem igualitária garante a mesma representatividade de todos os estratos na amostra, corrigindo possíveis grupos que porventura estejam subrepresentados.
    No entanto, caso a proporção de um estrato na população seja inferior à quantidade necessária, incorre na inviabilização da realização da amostragem. 
    \item Ela precisaria se preocupar menos com a construção do algoritmo computacional para a realização da amostragem por se tratar de um método mais complexo.
     Por outro lado, por se tratar de um método puramente aleatório, poderia ocorrer de selecionar apenas uma parcela específica da população que 
     poderia não ser representativa dentro do objetivo do estudo.  
\end{enumerate}

5) Considere os seguintes cenários:

\textbf{Cenário 1:}
Gabriel está realizando seu TCC e tem como objetivo verificar se os alunos do curso de Administração possuem conhecimento acerca do Mercado Financeiro. Para isso, ele cria 
um formulário utilizando a ferramenta Google Forms e o divulga em um grupo do Whatsapp dos alunos que entraram junto com ele na graduação. Após analisar as respostas do formulário, 
verifica que 70\% dos alunos responderam que conhecem o Mercado Financeiro muito bem e inclusive investem mensalmente em ações. 

\textbf{Cenário 2:}
Amanda participa da Liga de Investimentos da UFRJ. Foi passada para ela a tarefa de verificar se os alunos do curso de Administração possuem conhecimento acerca do Mercado Financeiro. 
O prazo para ela realizar essa tarefa é de no máximo 3 semanas. Amanda decide planejar sua pesquisa e se dá conta que seria impossível entrevistar 
todos os alunos do curso de Administração nesse intervalo de tempo, então ela estabelece a seguinte estratégia:

\begin{itemize}
    \item Pedir a cordenadora do curso de Administração a pauta contendo os alunos inscritos em cada uma das disciplinas do semestre atual conjuntamente com os e-mails cadastrados no SIGA.
    \item Colocar o nome de todos os alunos em uma tabela de excel (removendo duplicidades).
    \item Selecionar aleatoriamente 50 estudantes dessa lista (sem reposição).
    \item Enviar um e-mail convidando cada um dos 50 estudantes a responder um formulario criado através da ferramenta Google Forms. 
\end{itemize}

Após o período de 3 semanas, Amanda nota que 47 estudantes responderam seu questionário,
verificando que aproximadamente $34\%$ dos estudantes responderam que conhecem o Mercado Financeiro muito bem, inclusive investindo mensalmente em ações.   

Analisando os dois cenários, responda as seguintes perguntas:

\begin{enumerate}[a)]
    \item No Cenário 1, Gabriel pode ter um resultado viesado em sua pesquisa, haja vista que, como está realizando o TCC, provavelmente seus colegas de turma estão 
    em períodos próximos ou iguais ao último, tendo tido contato com matérias como Análise de Investimentos e Gerência Financeira. Dessa forma, Gabriel deixa de fora
    estudantes de períodos iniciais, criando um recorte da população que não condiz com a população objetivada no estudo. 
    
    No Cenário 2, por se basear em um método de amostragem probabilístico, não há como atribuir culpa a Amanda em relação à condução da pesquisa. Um apontamento
    que pode ser feito é que ela poderia ter estratificado a população com base nos diferentes períodos da formação, obtendo, dessa forma, respostas de estudantes 
    com níveis de formação distintos.  

    \item Não é possível afirmar categoricamente que o resultado encontrado por Amanda seja melhor do que o encontrado por Gabriel. A diferença é que podemos criticar
    Gabriel por sua escolha de pessoas para compor a amostra. Por outro lado, não podemos criticar Amanda por sua seleção, haja vista que não teve nenhuma influência sobre 
    as pessoas selecionadas. 
    \item Amostragem por Conveniência.
    \item Amostragem Aleatória Simples.
    \item Amostragem Aleatória Estratificada. 
\end{enumerate}

6) 

 \textbf{Mostrando que sempre conseguimos o tamanho amostral quando $k\frac{N}{n}$}

\vspace{3px}

 Na Amostragem Aleatória Sistemática, precisamos lembrar que, fixado o $k$,  um dos passos é sortear um inteiro aleatório entre $1$ e $k$ ao qual vamos chamar de $x$. 

 Para que consigamos obter o tamanho da amostra desejado, precisamos garantir que, ao realizar a seleção dos elementos da amostra, não iremos extrapolar o tamanho da população, 
 em outras palavras, basta verificar que:

 $$x + (n-1) \cdot k \leq N.$$

 Como  $k = \dfrac{N}{n}$, temos que a expressão acima pode ser reescrita da seguinte forma:

 $$x + (n-1) \cdot \dfrac{N}{n} \leq N.$$

 A partir disso:

  $$x + (n-1) \cdot \dfrac{N}{n} \leq N \implies x + N -\dfrac{N}{n} \leq N \implies x -\dfrac{N}{n} \leq 0 .$$

Sabemos que $x \leq k$, pois $x$ necessariamente é um número inteiro entre $1$ e $k$, logo $x -\dfrac{N}{n} \leq 0$ é sempre verdade e, portanto, a garantia está estabelecida. 


 \textbf{Mostrando que  quando  $k = \dfrac{N}{0,5 \cdot n}$ podemos incorrer em um tamanho amostral inferior ao desejado}

 Vamos realizar a mesmo procedimento do caso anterior, ou seja, temos um $x$ entre 1 e $k$ e vamos verificar o que acontece. Novamente, para que consigamos 
 selecionar a quantidade de elementos desejados, precisamos garantir que:

  $$x + (n-1) \cdot k \leq N.$$

Nesse caso, substituindo o valor de $k$ temos que:

  $$ x + (n-1) \cdot \dfrac{N}{0,5 \cdot n} \leq N.$$

Assim:

  $$ x + (n-1) \cdot \dfrac{N}{0,5 \cdot n} \leq N \implies  x + \dfrac{N}{0,5} - \dfrac{N}{0,5 \cdot n} \leq N \implies x - \dfrac{2N}{n} \leq -N$$

Nesse caso, temos que a desigualdade final não é sempre verdadeira, um caso particular que a invalida é tomar $x=k$, nesse caso teríamos que $0<-N$ o que é 
absrudo, dado que em um contexto prático $N$ sempre será maior que 0. 

Um exemplo prático pode ser observado no seguinte caso:

Suponha que desejamos uma amostra contendo $10$ idades de indivíduos de uma certa população de tamanho $100$ por meio da Amostragem Aleatória Sistmátoca.
 Definindo $k=\frac{N}{0,5 \cdot n}$ temos que $k=\dfrac{100}{5} = 20$

Vamos sortear um número aleatório entre $1$ e $20$, suponha que seja 12. Dessa forma, o elemento que ocupa a posição 12 na população será selecionado para minha amostra,
 e assim sucessivamente: 12, 22, 32, 42, 52, 62, 72, 82, 92, \textcolor{red}{opa!} se eu andar mais 10 casas, então chego no valor $102$, que ultrapassa o tamanho da população. 
 No fim, acabei ficando com uma amostra contendo 9 elementos, um tamanho menor do que eu desejava inicialmente. 


\vspace{3px}



\end{document}